\documentclass[12pt,oneside]{book}

% ============================================================
% METADATOS
% ============================================================
\newcommand{\universidad}{Universidad de Buenos Aires (UBA)}
\newcommand{\materia}{Derecho Civil --- Introducción y Parte General}
\newcommand{\titulo}{Fin de la Existencia de la Persona Humana}
\newcommand{\autor}{Isaac Edgar Camacho Ocampo}
\newcommand{\anio}{2025}

% ============================================================
% PAQUETES
% ============================================================
\usepackage[utf8]{inputenc}
\usepackage[T1]{fontenc}
\usepackage[spanish,es-nodecimaldot]{babel}

\usepackage[
    top=2.5cm,
    bottom=2.5cm,
    left=3cm,
    right=2.5cm,
    headheight=15pt
]{geometry}

\usepackage{amsmath}
\usepackage{amssymb}
\usepackage{mathtools}

\usepackage{graphicx}
\usepackage{float}
\usepackage{caption}
\usepackage{subcaption}

\usepackage{booktabs}
\usepackage{array}
\usepackage{longtable}
\usepackage{tabularx}

\usepackage{enumitem}

\usepackage{xcolor}
\usepackage[most]{tcolorbox}

\usepackage{fancyhdr}
\usepackage{ragged2e}

\usepackage[
    colorlinks=true,
    linkcolor=blue!50!black,
    citecolor=green!40!black,
    urlcolor=blue!60!black,
    pdftitle={\titulo},
    pdfauthor={\autor},
    pdfsubject={Apuntes universitarios},
    bookmarks=true
]{hyperref}

% ============================================================
% RUTAS DE IMÁGENES
% ============================================================
\graphicspath{
    {./img/}
    {./graficos/}
    {./figuras/}
}

% ============================================================
% CONFIGURACIÓN GENERAL
% ============================================================
\setcounter{tocdepth}{2}

\setlength{\parindent}{0pt}
\setlength{\parskip}{0.5em}

\setlist[itemize]{
    topsep=0.3em,
    itemsep=0.2em
}

\setlist[enumerate]{
    topsep=0.3em,
    itemsep=0.2em
}

\clubpenalty=10000
\widowpenalty=10000

% ============================================================
% ENCABEZADOS Y PIES
% ============================================================
\pagestyle{fancy}
\fancyhf{}

\fancyhead[L]{\nouppercase{\leftmark}}
\fancyhead[R]{\materia}

\fancyfoot[C]{\thepage}

\renewcommand{\headrulewidth}{0.4pt}
\renewcommand{\footrulewidth}{0pt}

\fancypagestyle{plain}{
    \fancyhf{}
    \fancyfoot[C]{\thepage}
    \renewcommand{\headrulewidth}{0pt}
}

% ============================================================
% COMANDOS MATEMÁTICOS
% ============================================================
\newcommand{\abs}[1]{\left|#1\right|}
\newcommand{\norm}[1]{\left\lVert#1\right\rVert}
\newcommand{\vect}[1]{\mathbf{#1}}
\newcommand{\deriv}[2]{\frac{d#1}{d#2}}
\newcommand{\pderiv}[2]{\frac{\partial#1}{\partial#2}}

% ============================================================
% CAJAS ACADÉMICAS
% ============================================================
\newtcolorbox{definition}[1]{
    enhanced,
    breakable,
    title={Definición: #1},
    fonttitle=\bfseries,
    colback=blue!3,
    colframe=blue!50!black,
    boxrule=0.8pt,
    arc=2mm
}

\newtcolorbox{important}{
    enhanced,
    breakable,
    title={Importante},
    fonttitle=\bfseries,
    colback=yellow!8,
    colframe=orange!70!black,
    boxrule=0.8pt,
    arc=2mm
}

\newtcolorbox{example}[1]{
    enhanced,
    breakable,
    title={Ejemplo: #1},
    fonttitle=\bfseries,
    colback=green!4,
    colframe=green!40!black,
    boxrule=0.8pt,
    arc=2mm
}

\newtcolorbox{formula}[1]{
    enhanced,
    breakable,
    title={Fórmula / Regla: #1},
    fonttitle=\bfseries,
    colback=purple!4,
    colframe=purple!50!black,
    boxrule=0.8pt,
    arc=2mm
}

\newtcolorbox{exercise}[1]{
    enhanced,
    breakable,
    title={Ejercicio: #1},
    fonttitle=\bfseries,
    colback=gray!5,
    colframe=gray!60!black,
    boxrule=0.8pt,
    arc=2mm
}

\newtcolorbox{note}{
    enhanced,
    breakable,
    title={Nota},
    fonttitle=\bfseries,
    colback=white,
    colframe=gray!50,
    boxrule=0.6pt,
    arc=2mm
}

% ============================================================
% DOCUMENTO
% ============================================================
\begin{document}

% ------------------------------------------------------------
% PORTADA
% ------------------------------------------------------------
\begin{titlepage}

\thispagestyle{empty}

\begin{center}

\vspace*{1cm}

{\large \textsc{\universidad}}

\vspace{3cm}

{\Huge \textbf{\materia}}

\vspace{1cm}

{\LARGE \titulo}

\vspace{0.8cm}

\textit{Comprender cómo el derecho define y comprueba la muerte es comprender los cimientos de todo el sistema sucesorio.}

\vspace{1.2cm}

\rule{0.70\textwidth}{0.5pt}

\vspace{0.5cm}

{\large Apuntes de estudio}

\vspace{0.5cm}

\rule{0.70\textwidth}{0.5pt}

\vfill

{\large \autor}

\vspace{0.3cm}

{\large \anio}

\end{center}

\vfill

\begin{flushleft}
\textit{Este es un modesto aporte a los estudiantes de ``Derecho Civil'' no pretende sustituir el material oficial de la materia.}
\end{flushleft}

\end{titlepage}

% ------------------------------------------------------------
% ÍNDICE
% ------------------------------------------------------------
\tableofcontents

% ------------------------------------------------------------
% CONTENIDO
% ------------------------------------------------------------
\chapter{Fin de la Existencia de la Persona Humana}

\section*{Ubicación temática y relevancia práctica}
\addcontentsline{toc}{section}{Ubicación temática y relevancia práctica}

La regulación civil de la persona humana comienza con el primer momento de su existencia --que en el ordenamiento argentino, en consonancia con las normas constitucionales, se sitúa en el momento de la concepción-- y concluye con el fin de la existencia, es decir, con la muerte. El estudio de esta última etapa presenta dos grandes cuestiones: por un lado, la situación de la muerte propiamente dicha y su acreditación civil; por otro, el instituto de la presunción de fallecimiento, figura jurídica aplicable cuando una persona se encuentra ausente de su domicilio sin que se tengan noticias de ella durante el plazo fijado por la ley. En la presente unidad se desarrolla exclusivamente el fin de la existencia por muerte comprobada; la presunción de fallecimiento corresponde a una unidad temática distinta.

\begin{note}
Los distintos bloques de esta unidad se encuentran conectados entre sí. Ante el fallecimiento de una persona es necesario, en primer lugar, comprender qué es la muerte y cómo se comprueba (Bloque I). Cuando no se dispone de un cadáver, corresponde analizar cómo se acredita judicialmente el fallecimiento (Bloque III). Si dos o más personas fallecen simultáneamente, la conmoriencia determina las relaciones sucesorias entre ellas (Bloque III). Finalmente, la muerte produce consecuencias jurídicas concretas sobre el patrimonio y los derechos de la persona fallecida (Bloque IV).
\end{note}

El fin de la existencia de la persona humana constituye uno de los pilares del derecho civil: el sistema de herencias, la donación de órganos, los registros civiles, los seguros de vida y las sucesiones dependen de determinar cuándo, cómo y con qué efectos jurídicos se produce la muerte. En el ejercicio profesional del derecho privado resulta indispensable gestionar partidas de defunción, iniciar sucesiones, asesorar a las familias en situaciones de presunción de fallecimiento, promover o impugnar declaraciones judiciales de muerte y conocer los límites de la transmisión hereditaria.

\section{La muerte: concepto, comprobación y normativa}

\subsection{Principio general}

El Código Civil y Comercial de la Nación (en adelante, CCyCN) establece el principio general que rige la finalización de la existencia de la persona humana:

\begin{formula}{Art. 93 CCyCN --- Principio general}
La existencia de la persona humana termina por su muerte.
\end{formula}

El Código Civil anterior contenía una disposición similar, aunque hacía referencia a la ``muerte natural''. El término ``natural'', en la redacción del Código de Vélez Sarsfield, no se oponía a ``provocada'' --como si una persona asesinada no muriera-- sino que aludía a la muerte biológica y, fundamentalmente, procuraba excluir cualquier forma de muerte civil.

\begin{note}
La muerte civil implicaba que el derecho consideraba muerta a una persona por el solo cumplimiento de determinada circunstancia --por ejemplo, una pena, o el ingreso a una comunidad religiosa, especialmente de votos solemnes, situación en la cual la persona ``moría al mundo'' para entrar en religión--. Esta institución nunca tuvo aplicación efectiva en el derecho argentino desde el Código Civil anterior, y el CCyCN simplifica la cuestión disponiendo, sin más, que la existencia de la persona termina con su muerte.
\end{note}

\subsection{¿Qué es la muerte?}

\begin{definition}{Muerte}
Cese irreversible de las funciones corporales; fin de la unidad de cuerpo y alma. El momento exacto en que la muerte ocurre permanece, en gran medida, en el plano del misterio: el derecho civil no define en qué instante preciso se produce, sino que regula el modo en que se \textbf{comprueba} su acaecimiento, comprobación que siempre resulta posterior al hecho biológico en sí.
\end{definition}

El momento exacto de la muerte de una persona resulta, por lo tanto, muy difícil de determinar. Aquello de lo que se ocupa el derecho civil es de su comprobación.

\subsection{Comprobación de la muerte}

\begin{formula}{Art. 94 CCyCN --- Comprobación de la muerte}
La comprobación de la muerte queda sujeta a los estándares médicos aceptados, aplicándose la legislación especial en el caso de ablación de órganos del cadáver.
\end{formula}

La comprobación de la muerte se realiza conforme a estándares médicos que van variando a lo largo del tiempo, a medida que avanza la ciencia, y que permiten detectar el cese irreversible de las funciones corporales --especialmente las cardiorrespiratorias y cerebrales, que son las que comandan el organismo humano.

Se habla de ``comprobación'' precisamente porque esta ocurre con posterioridad al hecho de la muerte. Lo que se obtiene es una certeza moral, no una certeza apodíctica: las certezas sobre la muerte se basan siempre en observaciones empíricas que, en principio, podrían tener alguna falla, pero que --con el acompañamiento de los estándares médicos-- se procuran hacer lo más rigurosas posible, de manera de no dejar duda alguna acerca del fallecimiento de la persona. De allí la confianza depositada en la intervención médica.

\subsection{Instrumentos para acreditar la muerte}

\begin{formula}{Ley 26.413 --- Registro del Estado Civil y Capacidad de las Personas}
Regula la forma en que deben ser expedidos el certificado médico de defunción y la partida de defunción. El certificado médico es expedido por un profesional de la salud siempre sobre la base de la presencia del cadáver. La partida de defunción es el instrumento público en el que consta el asiento en el registro civil, con la fecha, la hora y los demás elementos que rodean a la muerte.
\end{formula}

La normativa distingue con claridad dos instrumentos:

\begin{table}[H]
\centering
\caption{Certificado médico de defunción y partida de defunción}
\begin{tabularx}{\textwidth}{X X}
\toprule
\textbf{Certificado médico de defunción} & \textbf{Partida de defunción} \\
\midrule
Expedido por un profesional de la salud & Instrumento público del Registro Civil \\
Siempre requiere la presencia del cadáver & No requiere la presencia del cadáver (se basa en el certificado) \\
Documento de base para inscribir la muerte & Constancia del asiento en el Registro del Estado Civil \\
Contiene fecha, hora y elementos que rodean la muerte & Regulada también por la Ley 26.413 \\
\bottomrule
\end{tabularx}
\end{table}

\section{Donación de órganos y criterios de muerte encefálica}

\subsection{Consentimiento presunto}

Cuando una persona fallece sin haber expresado su voluntad respecto de la donación de órganos, la ley presume que es dadora.

\begin{formula}{Ley 27.447 (2018) --- Trasplante de Órganos y Tejidos}
Establece que cuando una persona fallece sin haber expresado su voluntad, se presume que es dadora de órganos. La negativa debe ser expresa, realizada ante los organismos señalados en la ley. El artículo 36 dispone que la muerte puede certificarse con la confirmación del cese irreversible de las funciones circulatorias o encefálicas. El artículo 37 remite a un protocolo fijado por el Ministerio de Salud con asesoramiento del INCUCAI (Resolución 716/2019), estableciendo que la certificación debe ser realizada por dos médicos, uno de los cuales debe ser neurólogo o neurocirujano y no puede integrar el equipo de ablación.
\end{formula}

\begin{note}
La Ley 27.447 presume que toda persona es donante de órganos, salvo que haya expresado directamente una negativa ante los organismos señalados en la norma. Esta solución ha generado controversias, tanto por suprimir el rol de la familia en la decisión como por presumir una voluntad que, para algunos, debería quedar estimulada en lugar de presumida.
\end{note}

\subsection{Antecedentes legislativos}

\begin{formula}{Ley 21.541 (1977) --- Primera ley de trasplantes}
Definió el fallecimiento como la disfunción total e irreversible de las funciones encefálicas, verificada por un equipo médico. La reglamentación establecía los signos mínimos que debían presentarse en su totalidad.
\end{formula}

\begin{formula}{Ley 24.193 (1993) --- Ley de trasplantes anterior a 2018}
Establecía que la muerte se constataba por: ausencia irreversible de respuesta cerebral con pérdida absoluta de conciencia; ausencia de respiración espontánea; ausencia de reflejos cefálicos; pupilas fijas no reactivas; e inactividad encefálica comprobada por medios técnicos o instrumentales (electroencefalograma plano y constante durante seis horas). No era necesario comprobar la inactividad encefálica en caso de paro cardiorrespiratorio total e irreversible.
\end{formula}

Esta definición de la Ley 24.193 estuvo vigente entre 1993 y 2018.

\subsection{Criterios actuales y Protocolo 2019}

A partir de 2018, la ley opta por no incorporar en su texto una descripción tan detallada de los signos, remitiéndola a un protocolo. Se dispone que el fallecimiento de una persona puede certificarse con la confirmación del cese irreversible de las funciones circulatorias o encefálicas --cualquiera de los dos supuestos, tanto el paro cardiorrespiratorio como el encefálico.

\begin{important}
La palabra clave del criterio vigente es la \textbf{irreversibilidad}. Si existe la posibilidad de reversión --como ocurre con un paro cardíaco que puede revertirse mediante reanimación-- no hay persona fallecida, sino una persona cuya función circulatoria puede ser terapéuticamente reanimada. El fallecimiento, en sentido propio, es un cese irreversible de las funciones.
\end{important}

Respecto de la certificación, el artículo 37 de la Ley 27.447 remite a un protocolo fijado por el Ministerio de Salud con asesoramiento del INCUCAI, concretado en la Resolución 716/2019. Dicho protocolo puede variar, en atención a la necesidad de actualización técnica de los medios biotecnológicos utilizados para el diagnóstico.

\subsection{Garantías y críticas al régimen vigente}

\begin{important}
Un sector de la doctrina ha criticado el sistema vigente por no establecer en el propio texto legal algunas garantías mínimas --como las que preveía la ley anterior--, delegando en una reglamentación un hecho tan decisivo como la certificación de la muerte. Esta cuestión ha merecido discusión doctrinaria; no obstante, la ley mantiene como criterio decisivo el del cese irreversible de las funciones.
\end{important}

El desarrollo del apunte remite, para la profundización de este debate, a los trabajos del Dr.\ José Tobías y del Dr.\ Carlos Muñiz, quienes han escrito ampliamente sobre la determinación del momento de la muerte.

\section{Conmoriencia y prueba supletoria del fallecimiento}

\subsection{Conmoriencia}

\begin{formula}{Art. 95 CCyCN --- Conmoriencia}
Se presume que las personas que fallecen en un desastre común mueren al mismo tiempo. Salvo prueba en contrario, es una presunción iuris tantum.
\end{formula}

\begin{note}
A lo largo del tiempo han existido distintas tesis sobre la preeminencia --es decir, sobre quién moría primero, si los más jóvenes, los más ancianos, o de acuerdo con el sexo--. Los códigos civiles argentinos, ya desde la época de Vélez Sarsfield, establecen que todas las personas que se encuentran en esas circunstancias mueren al mismo tiempo, admitiéndose prueba en contrario.
\end{note}

\begin{example}{Accidente de tránsito}
Fallecen el padre y el hijo en un accidente de tránsito; sobrevive la madre. Se supone que no existen otros herederos.

Si el hijo fallece \textbf{después} del padre: hereda del padre y, al morir el propio hijo, lo hereda la madre.

Si el hijo fallece \textbf{antes} que el padre: el padre, al morir, no tendría herederos directos, y serían sus propios ascendientes quienes heredarían.

\textbf{Solución legal (Art. 95 CCyCN):} se presume que ambos fallecieron al mismo tiempo. No hay transmisión de derechos entre ellos, lo cual evita discusiones sobre quién murió primero y alteraciones en los órdenes sucesorios.
\end{example}

\subsection{Prueba supletoria del fallecimiento}

El segundo supuesto especial que corresponde analizar es el previsto en el artículo 98, segunda parte, del CCyCN, que no debe confundirse con la presunción de fallecimiento.

\begin{formula}{Art. 98 CCyCN --- Prueba supletoria del fallecimiento}
Cuando el cadáver no es hallado o no puede ser identificado, y la desaparición se produjo en circunstancias tales que la muerte debe tenerse por cierta, el juez puede tener por comprobada la muerte y disponer la correspondiente inscripción en el registro.
\end{formula}

Para inscribir una muerte se requiere, en principio, un certificado médico. Sin embargo, cuando el profesional de la salud no dispone de cadáver alguno, o no puede identificarlo, dicha vía resulta imposible. Ante la certeza de que la persona ha fallecido, se habilita entonces un trámite judicial.

El criterio decisivo para optar por este artículo, en lugar de la presunción de fallecimiento, es la certeza:

\begin{table}[H]
\centering
\caption{Prueba supletoria (Art. 98) frente a la presunción de fallecimiento}
\begin{tabularx}{\textwidth}{X X}
\toprule
\textbf{Prueba supletoria --- Art. 98} & \textbf{Presunción de fallecimiento} \\
\midrule
Hay \textbf{certeza} de que la persona falleció & Hay \textbf{duda o incertidumbre} sobre si falleció \\
Cadáver no hallado o no identificado & Persona ausente sin noticias durante el tiempo fijado por ley \\
Procedimiento judicial más breve & Procedimiento judicial mucho más largo \\
El juez ordena la inscripción de la muerte & Se declara presuntamente fallecida a la persona \\
\bottomrule
\end{tabularx}
\end{table}

\subsection{Ejemplos desarrollados}

\begin{example}{Avión que se desintegra en el aire (Art. 98)}
Un avión de Air France que se desintegró en el aire (vuelo AF 447), o el avión de LAPA en el Río Uruguay que se desintegró en el aire: no hubo sobrevivientes, y resulta imposible que los hubiera habido. Ante ese hecho, se aplica el artículo 98 --prueba supletoria-- porque existe certeza de que las personas fallecieron.
\end{example}

\begin{example}{Los Andes (presunción de fallecimiento)}
El caso del avión de la película ``Viven'', cuyos pasajeros --uruguayos-- caen en la cordillera de los Andes: no solo no se los encontraba, sino que tampoco se sabía si estaban muertos o no. No podía afirmarse que esas personas habían fallecido; correspondía, por lo tanto, aplicar la presunción de fallecimiento y esperar los plazos propios de ese instituto. De hecho, fueron encontrados a los tres meses.
\end{example}

\begin{example}{Alpinista en precipicio inaccesible (Art. 98)}
Un alpinista cae en un precipicio de acceso imposible, acompañado por compañeros de montañismo que certifican que la persona llegó a ese punto, tuvo un problema, se resbaló y cayó, resultando imposible que haya sobrevivido e imposible acceder al lugar. Con esos testimonios y las constancias del viaje puede realizarse la prueba supletoria: no corresponde permanecer en la incertidumbre, pues la persona falleció. El juez, teniendo por comprobada la muerte, dispone la inscripción en el registro.
\end{example}

\section{Efectos jurídicos de la muerte}

\subsection{Efectos extrapatrimoniales}

La muerte produce, en el campo extrapatrimonial, la extinción de los atributos de la persona: el nombre y la capacidad. Del domicilio subsiste el domicilio convencional --aquel que se elige para un contrato-- que se transmite a los sucesores conforme al artículo 1024 CCyCN. Se extingue el estado civil, pero subsisten los vínculos a los fines sucesorios: si existía cónyuge, y respecto del padre, la madre y los hijos, el parentesco se mantiene a esos efectos.

\subsection{Protección jurídica post mortem}

Se extinguen los derechos personalísimos, pero subsiste una protección jurídica sobre el honor de la persona fallecida, su intimidad y la inviolabilidad de sus papeles privados. Existe, asimismo, una protección jurídica del cadáver, orientada a que este sea tratado con dignidad.

\begin{note}
El cadáver posee un estatuto especial: no es la persona ni es el cuerpo, pero constituye una prolongación de la persona, y por ello existe un cuidado de esta y de su memoria. La protección de la memoria de los difuntos se vincula también con su intimidad y con la inviolabilidad de sus papeles privados.
\end{note}

\subsection{Efectos patrimoniales}

En el campo patrimonial, los derechos patrimoniales de la persona fallecida se transmiten a sus herederos conforme a las reglas dispuestas por el CCyCN, que pueden ser de dos tipos: testamentarias --cuando la propia persona ha determinado quiénes son sus herederos, existiendo porciones legítimas correspondientes a herederos forzosos sobre las cuales no puede disponerse, y una parte disponible según existan descendientes o solo ascendientes-- o, en su defecto, ab intestato.

\subsection{Derechos no transmisibles}

\begin{formula}{Art. 1024 CCyCN --- Efectos de los contratos: sucesores}
Las obligaciones que son inherentes a las personas no son transmisibles por causa de muerte.
\end{formula}

\begin{formula}{Art. 2280 CCyCN --- Situación de los herederos}
No se transmiten por sucesión los derechos que la ley dispone que no sean transmisibles.
\end{formula}

No se transmiten, por tanto, los derechos que la ley dispone como no transmisibles. Un ejemplo de ello son las obligaciones inherentes a las personas, que no son transmisibles por causa de muerte.

\begin{example}{El pintor famoso}
Se contrata a un pintor famoso para que realice el retrato de un familiar por cinco mil dólares. El pintor fallece. Su hija, también pintora, ofrece realizar el retrato en su lugar. La obligación no se transmite, porque era inherente a la persona --se habían considerado las cualidades específicas del pintor contratado-- y su transmisión resultaría incompatible con la naturaleza de la obligación pactada.
\end{example}

Los derechos reales de usufructo y de habitación constituyen otros ejemplos típicos de derechos que no se transmiten por causa de muerte: por su propia constitución, se extinguen con la muerte del beneficiario. La muerte produce, por regla general, la transmisión de los derechos patrimoniales, con excepción de aquellos que no se transmiten.

\section*{Cierre y síntesis de la clase}
\addcontentsline{toc}{section}{Cierre y síntesis de la clase}

La unidad ha desarrollado las cuestiones relativas al fin de la existencia de la persona humana: la muerte, su comprobación, el caso de la conmoriencia, la prueba supletoria del fallecimiento y los efectos jurídicos de la muerte.

% ------------------------------------------------------------
% CUADRO CORNELL
% ------------------------------------------------------------
\section*{Cuadro Cornell --- Repaso de la unidad}
\addcontentsline{toc}{section}{Cuadro Cornell --- Repaso de la unidad}

\begin{longtable}{
    >{\raggedright\arraybackslash}p{0.30\textwidth}
    >{\raggedright\arraybackslash}p{0.60\textwidth}
}
\caption{Cuadro Cornell --- Fin de la existencia de la persona humana} \\
\toprule
\textbf{Preguntas / palabras clave} & \textbf{Notas / respuestas} \\
\midrule
\endfirsthead

\toprule
\textbf{Preguntas / palabras clave} & \textbf{Notas / respuestas} \\
\midrule
\endhead

\textbf{¿Cuándo termina la existencia de la persona? (Art. 93)} &
Con la muerte. El CCyCN simplifica la redacción anterior eliminando la referencia a la ``muerte natural'' y la figura de la muerte civil, ya inexistente en el derecho argentino desde el Código de Vélez. \\

\textbf{¿Qué es la muerte?} &
El cese irreversible de las funciones corporales. El momento exacto en que ocurre permanece en el plano del misterio; lo que el derecho regula es su comprobación posterior mediante estándares médicos. \\

\textbf{Art. 94 --- ¿Cómo se comprueba la muerte?} &
La comprobación queda sujeta a los estándares médicos aceptados, que evolucionan con la ciencia. Para la ablación de órganos rige la legislación especial (Ley 27.447 y Protocolo 2019). \\

\textbf{Certificado médico vs.\ partida de defunción} &
El certificado médico es expedido por un profesional de la salud, siempre ante la presencia del cadáver, y constituye el documento base. La partida de defunción es el instrumento público del Registro Civil donde consta la inscripción de la muerte (Ley 26.413). \\

\textbf{Ley 27.447 --- ¿Todos somos donantes de órganos?} &
Sí, por presunción legal. Quien no quiera serlo debe expresar una negativa ante los organismos designados. Se eliminó el rol de la familia que existía bajo la Ley 24.193 (Arts.\ 36 y 37). \\

\textbf{¿Qué exigía la Ley 24.193 (1993--2018) para la muerte encefálica?} &
Ausencia irreversible de respuesta cerebral; ausencia de respiración espontánea; ausencia de reflejos cefálicos; pupilas fijas no reactivas; encefalograma plano y constante durante seis horas. El paro cardiorrespiratorio total e irreversible eximía de la prueba encefálica. \\

\textbf{Ley 27.447 --- ¿cuál es el criterio actual de muerte?} &
Cese irreversible de las funciones circulatorias o encefálicas. Los signos específicos y los plazos se remiten al Protocolo del Ministerio de Salud (Res.\ 716/2019), elaborado con el INCUCAI. Debe certificarlo un equipo de dos médicos, uno de ellos neurólogo o neurocirujano ajeno al equipo de ablación. \\

\textbf{¿Qué es la conmoriencia? (Art. 95)} &
Presunción iuris tantum de que dos personas que fallecen en un desastre común mueren al mismo tiempo. Evita discusiones sobre quién murió primero y posibles alteraciones en los órdenes sucesorios; admite prueba en contrario. \\

\textbf{Prueba supletoria del fallecimiento (Art. 98, segunda parte)} &
Procede cuando el cadáver no es hallado o no puede identificarse, pero existe certeza de que la persona falleció. Un juez --no un médico-- ordena la inscripción de la muerte en el registro. Ejemplos: aviones desintegrados, alpinistas en precipicios inaccesibles. \\

\textbf{Diferencia entre el Art. 98 y la presunción de fallecimiento} &
Art.\ 98: hay certeza de la muerte, el procedimiento es más breve y el juez ordena la inscripción. Presunción de fallecimiento: hay duda o incertidumbre, el procedimiento es mucho más largo y se declara ``presuntamente fallecida'' a la persona. \\

\textbf{Efectos extrapatrimoniales de la muerte} &
Se extinguen el nombre, la capacidad, el estado civil y los derechos personalísimos. Subsisten el domicilio convencional (para contratos), los vínculos familiares a fines sucesorios y la protección jurídica del honor, la intimidad y la inviolabilidad del cadáver. \\

\textbf{Efectos patrimoniales de la muerte} &
Se transmiten a los herederos los derechos patrimoniales, conforme a las reglas testamentarias o ab intestato del CCyCN, respetando las porciones legítimas de los herederos forzosos. \\

\textbf{Derechos NO transmisibles por causa de muerte} &
Las obligaciones inherentes a la persona (intuitu personae) y los derechos reales de usufructo y de habitación (Arts.\ 1024 y 2280 CCyCN). Ejemplo: si se contrata a un pintor famoso y este fallece, su hija no está obligada a cumplir ni tiene derecho a exigir el pago. \\

\midrule
\multicolumn{2}{p{\dimexpr\textwidth-2\tabcolsep}}{
\textbf{Resumen:} El fin de la existencia de la persona humana se produce con la muerte, cuya única definición jurídica en el CCyCN es el cese irreversible de las funciones corporales (Art.\ 93). Como ese momento exacto es de difícil determinación, el derecho regula su comprobación: en el caso ordinario, mediante el certificado médico de defunción --que exige la presencia del cadáver-- y la posterior partida de defunción en el Registro Civil, ambos regulados por la Ley 26.413. Para la ablación de órganos, los criterios son más rigurosos y están regulados por la Ley 27.447 (2018) y el Protocolo del INCUCAI de 2019: basta el cese irreversible de las funciones circulatorias o encefálicas, certificado por dos médicos, uno de los cuales debe ser neurólogo o neurocirujano y ajeno al equipo de trasplante. Bajo esta ley, toda persona es donante presunta salvo negativa expresa. Cuando dos personas fallecen en un desastre común sin poder determinarse quién murió primero, la conmoriencia (Art.\ 95) presume que ambas murieron simultáneamente, evitando distorsiones en los órdenes sucesorios. Si el cadáver no es hallado pero hay certeza de la muerte, el juez puede tenerla por acreditada y ordenar su inscripción mediante la prueba supletoria del Art.\ 98, que se distingue de la presunción de fallecimiento por la ausencia de incertidumbre. Finalmente, la muerte produce efectos extrapatrimoniales --extinción del nombre, la capacidad y el estado civil, con subsistencia de los vínculos familiares a fines sucesorios y protección jurídica post mortem del honor y la intimidad-- y patrimoniales: transmisión del patrimonio a los herederos según las reglas del CCyCN, con excepción de los derechos inherentes a la persona (usufructo, habitación, obligaciones intuitu personae), que se extinguen con la muerte del titular.
} \\
\bottomrule
\end{longtable}

\end{document}
